% =========================================================
%  ppt/07-finite-state-machine.tex
%  Materi 7 — Finite State Machine untuk NPC
%  Bab IX buku ajar
% =========================================================
\documentclass[aspectratio=169,10pt]{beamer}
% =========================================================
%  ppt/config/packages.tex
%  Paket bersama untuk seluruh presentasi kuliah
%  Pengantar Pemrograman Game
% =========================================================

\usepackage[T1]{fontenc}
\usepackage[utf8]{inputenc}
\usepackage[indonesian]{babel}
\usepackage{lmodern}
\usepackage{textcomp}
\usepackage{amsmath}
\usepackage{amssymb}
\usepackage{graphicx}
\usepackage{booktabs}
\usepackage{tabularx}
\usepackage{array}
\usepackage{multicol}
\usepackage{listings}
\usepackage{upquote}
\usepackage{adjustbox}
\usepackage{tcolorbox}
\usepackage{tikz}
\usetikzlibrary{
  arrows.meta,
  positioning,
  shapes.geometric,
  shapes.misc,
  calc,
  fit,
  backgrounds,
  matrix,
  chains,
  decorations.pathreplacing,
  patterns
}

\tcbuselibrary{skins}

% Jangan memuat geometry: ukuran halaman diatur oleh beamer.

% =========================================================
%  ppt/config/theme.tex
%  Satu tema Beamer untuk keempat belas presentasi
%  Pengantar Pemrograman Game
% =========================================================

% --------------------------------------------------------- warna
\definecolor{ubNavy}{HTML}{1F3864}
\definecolor{ubAccent}{HTML}{2E75B6}
\definecolor{ubHead}{HTML}{D9E2F3}
\definecolor{ubSoft}{HTML}{F2F6FC}
\definecolor{ubLine}{HTML}{BFBFBF}
\definecolor{ubText}{HTML}{1F2937}
\definecolor{ubMuted}{HTML}{5A6678}
\definecolor{ubGreen}{HTML}{1E7B45}
\definecolor{ubGreenBg}{HTML}{F1F8F2}
\definecolor{ubRed}{HTML}{B3261E}
\definecolor{ubRedBg}{HTML}{FDF2F1}
\definecolor{ubWarn}{HTML}{A35A05}
\definecolor{ubWarnBg}{HTML}{FFF8EC}
\definecolor{ubGray}{HTML}{F6F6F6}
\definecolor{ubCodeBg}{HTML}{10233F}
\definecolor{ubCodeText}{HTML}{E8EEF7}
\definecolor{ubCodeKey}{HTML}{8FBCF0}
\definecolor{ubCodeCom}{HTML}{8FB37A}
\definecolor{ubCodeStr}{HTML}{D9A46B}
\definecolor{ubCodeNum}{HTML}{7E99BE}

% Nama warna yang dipakai gambar TikZ buku ajar.
\definecolor{ubPrimary}{HTML}{1F3864}
\definecolor{ubAlt}{HTML}{F2F6FC}

% --------------------------------------------------------- dasar
\usetheme{default}
\usefonttheme{professionalfonts}
\setbeamertemplate{navigation symbols}{}
\setbeamerfont{frametitle}{size=\Large,series=\bfseries}
\setbeamerfont{framesubtitle}{size=\small}
\setbeamerfont{footline}{size=\tiny}
\setbeamerfont{block title}{size=\small,series=\bfseries}
\setbeamerfont{block body}{size=\small}
\setbeamerfont{itemize/enumerate body}{size=\small}
\setbeamerfont{itemize/enumerate subbody}{size=\footnotesize}

\setbeamersize{text margin left=0.60cm,text margin right=0.60cm}
\uchyph=0

\setbeamercolor{normal text}{fg=ubText,bg=white}
\setbeamercolor{structure}{fg=ubAccent}
\setbeamercolor{frametitle}{fg=ubNavy,bg=white}
\setbeamercolor{title}{fg=ubNavy}
\setbeamercolor{footline}{fg=ubMuted,bg=white}
\setbeamercolor{itemize item}{fg=ubAccent}
\setbeamercolor{itemize subitem}{fg=ubAccent}
\setbeamercolor{block title}{fg=white,bg=ubNavy}
\setbeamercolor{block body}{fg=ubText,bg=ubSoft}
\setbeamercolor{block title alerted}{fg=white,bg=ubWarn}
\setbeamercolor{block body alerted}{fg=ubText,bg=ubWarnBg}
\setbeamercolor{block title example}{fg=white,bg=ubGreen}
\setbeamercolor{block body example}{fg=ubText,bg=ubGreenBg}
\setbeamercolor{caption name}{fg=ubMuted}

\setbeamertemplate{blocks}[rounded][shadow=false]
\setbeamertemplate{itemize item}{\raisebox{0.10em}{\color{ubAccent}\scriptsize$\blacksquare$}}
\setbeamertemplate{itemize subitem}{\raisebox{0.06em}{\color{ubAccent}\tiny$\blacktriangleright$}}
\setbeamertemplate{enumerate item}{%
  \makebox[1.3em][l]{\color{ubAccent}\insertenumlabel.}}

\addtobeamertemplate{itemize/enumerate body begin}{%
  \setlength{\itemsep}{0.42em}\setlength{\parskip}{0pt}}{}

% --------------------------------------------------------- judul frame
\setbeamertemplate{frametitle}[default][left]

% --------------------------------------------------------- kaki
\newcommand{\materino}{00}
\newcommand{\materifoot}{Pengantar Pemrograman Game \textperiodcentered\ Materi \materino}

\setbeamertemplate{footline}{%
  \begin{beamercolorbox}[wd=\paperwidth,ht=0.015cm,dp=0cm]{footline}%
    \color{ubHead}\rule{\paperwidth}{0.4pt}%
  \end{beamercolorbox}%
  \begin{beamercolorbox}[wd=\paperwidth,ht=0.40cm,dp=0.14cm,%
                         leftskip=0.60cm,rightskip=0.60cm]{footline}%
    \usebeamerfont{footline}\color{ubMuted}\materifoot\hfill\insertframenumber%
  \end{beamercolorbox}%
}

% --------------------------------------------------------- makro buku ajar
\providecommand{\code}[1]{\texttt{\detokenize{#1}}}
\providecommand{\istilah}[1]{\emph{#1}}
\providecommand{\CPMK}[1]{CPMK\,#1}
\providecommand{\CPL}[1]{CPL\,#1}
\providecommand{\SUBCPMK}[1]{Sub-CPMK\,#1}

% --------------------------------------------------------- gaya kode
\lstdefinelanguage{CSharp}{
  morekeywords={abstract,as,base,bool,break,byte,case,catch,char,class,const,
    continue,decimal,default,delegate,do,double,else,enum,event,explicit,
    extern,false,finally,fixed,float,for,foreach,goto,if,implicit,in,int,
    interface,internal,is,lock,long,namespace,new,null,object,operator,out,
    override,params,private,protected,public,readonly,ref,return,sbyte,
    sealed,short,sizeof,stackalloc,static,string,struct,switch,this,throw,
    true,try,typeof,uint,ulong,unchecked,unsafe,ushort,using,virtual,void,
    volatile,while,var,get,set},
  sensitive=true,
  morecomment=[l]{//},
  morecomment=[s]{/*}{*/},
  morestring=[b]",
}

\lstdefinestyle{kode}{
  basicstyle=\ttfamily\footnotesize\color{ubCodeText},
  keywordstyle=\color{ubCodeKey}\bfseries,
  commentstyle=\color{ubCodeCom}\itshape,
  stringstyle=\color{ubCodeStr},
  numberstyle=\tiny\color{ubCodeNum},
  numbers=left,
  numbersep=7pt,
  backgroundcolor=\color{ubCodeBg},
  frame=none,
  showstringspaces=false,
  breaklines=false,
  tabsize=4,
  columns=fullflexible,
  keepspaces=true,
  xleftmargin=12pt,
  xrightmargin=2pt,
  aboveskip=0pt,
  belowskip=0pt,
  upquote=true,
}
\lstset{style=kode}

\newcommand{\kodejelas}[1]{%
  \vspace{0.16cm}%
  {\footnotesize\color{ubMuted}#1\par}%
}
\newcommand{\sumber}[1]{%
  \vspace{0.08cm}%
  {\scriptsize\color{ubMuted}\itshape #1\par}%
}
\newcommand{\awalan}[1]{%
  {\small\color{ubMuted}#1\par}\vspace{0.18cm}%
}

% --------------------------------------------------------- kotak
\newtcolorbox{kotak}[1][]{enhanced, colback=ubSoft, colframe=ubAccent,
  boxrule=0pt, leftrule=2.2pt, arc=1pt, left=6pt, right=6pt, top=4pt,
  bottom=4pt, fontupper=\footnotesize, #1}
\newtcolorbox{peringatan}[1][]{enhanced, colback=ubWarnBg, colframe=ubWarn,
  boxrule=0pt, leftrule=2.2pt, arc=1pt, left=6pt, right=6pt, top=4pt,
  bottom=4pt, fontupper=\footnotesize, #1}
\newtcolorbox{kiat}[1][]{enhanced, colback=ubGreenBg, colframe=ubGreen,
  boxrule=0pt, leftrule=2.2pt, arc=1pt, left=6pt, right=6pt, top=4pt,
  bottom=4pt, fontupper=\footnotesize, #1}

% --------------------------------------------------------- sampul
% \deckcover{nomor}{minggu}{judul}{sub-cpmk}{teknologi}{sumber}{blok}
\newcommand{\deckcover}[7]{%
\begin{frame}[plain]
\begin{tikzpicture}[remember picture, overlay]
  \fill[ubNavy] (current page.south west) rectangle ([xshift=6.00cm]current page.north west);
  \fill[ubAccent] ([xshift=6.00cm]current page.south west) rectangle ([xshift=6.08cm]current page.north west);

  \node[anchor=north west, text width=4.72cm, align=left, inner sep=0pt]
    at ($(current page.north west)+(0.66cm,-0.60cm)$) {%
      {\color{white}\bfseries\fontsize{10.2}{12.2}\selectfont
       Universitas Bale Bandung}\\[4pt]
      {\color{ubHead}\fontsize{7.6}{9.4}\selectfont
       Fakultas Teknologi Informasi}\\[1pt]
      {\color{ubHead}\fontsize{7.6}{9.4}\selectfont
       Program Studi Teknik Informatika}\\[7pt]
      {\color{ubAccent}\rule{4.0cm}{0.5pt}}%
    };

  \node[anchor=north west, text width=4.72cm, align=left, inner sep=0pt]
    at ($(current page.north west)+(0.66cm,-2.72cm)$) {%
      {\color{ubAccent}\fontsize{6.4}{7.8}\selectfont MATA KULIAH}\\[3pt]
      {\color{white}\bfseries\fontsize{9.8}{12.0}\selectfont
       Pengantar Pemrograman Game}\\[3pt]
      {\color{ubHead}\fontsize{7.6}{9.4}\selectfont
       MWCA009 \textperiodcentered\ 3 SKS}\\[8pt]
      {\color{ubHead}\fontsize{7.6}{9.4}\selectfont
       Semester Ganjil 2026/2027}%
    };

  \node[anchor=south west, text width=4.72cm, align=left, inner sep=0pt]
    at ($(current page.south west)+(0.66cm,0.52cm)$) {%
      {\color{ubAccent}\fontsize{6.4}{7.8}\selectfont DOSEN PENGAMPU}\\[3pt]
      {\color{white}\fontsize{8.4}{10.4}\selectfont
       Cecep Suwanda, S.Si., M.Kom.}%
    };

  \node[anchor=north west, text width=8.95cm, align=left, inner sep=0pt]
    at ($(current page.north east)+(-9.40cm,-0.92cm)$) {%
      {\color{ubAccent}\bfseries\fontsize{8.2}{10.2}\selectfont
       MATERI #1\ \ \textperiodcentered\ \ MINGGU #2}\\[8pt]
      {\color{ubNavy}\bfseries\fontsize{19.4}{23.4}\selectfont
       \hyphenpenalty=10000 \exhyphenpenalty=10000 \raggedright #3}%
    };

  \node[anchor=north west, text width=8.95cm, align=left, inner sep=0pt]
    at ($(current page.north east)+(-9.40cm,-4.30cm)$) {%
      \renewcommand{\arraystretch}{1.30}%
      \begin{tabular}{@{}>{\color{ubMuted}\fontsize{7.3}{9.0}\selectfont}p{1.85cm}
                        >{\color{ubText}\fontsize{8.1}{10.0}\selectfont}p{6.7cm}@{}}
        Sub-CPMK      & #4 \\
        Teknologi     & #5 \\
        Sumber materi & #6 \\
        Blok          & #7 \\
      \end{tabular}%
    };
\end{tikzpicture}
\end{frame}%
}

% --------------------------------------------------------- pembatas bagian
\newcommand{\sectionframe}[3]{%
\begin{frame}[plain]
\begin{tikzpicture}[remember picture, overlay]
  \fill[ubNavy] (current page.south west) rectangle (current page.north east);
  \fill[ubAccent] (current page.north west) rectangle ([xshift=0.10cm]current page.south west);
  \node[anchor=west, inner sep=0pt]
    at ($(current page.west)+(1.00cm,1.45cm)$)
    {\color{ubNavy!52!white}\bfseries\fontsize{58}{58}\selectfont #1};
  \node[anchor=west, inner sep=0pt, text width=11.6cm, align=left]
    at ($(current page.west)+(1.05cm,-0.55cm)$) {%
      {\color{ubAccent}\bfseries\fontsize{8.0}{10.0}\selectfont BAGIAN}\\[3pt]
      {\color{white}\bfseries\fontsize{20.5}{24.5}\selectfont #2}\\[5pt]
      {\color{ubHead}\small #3}%
    };
  \node[anchor=south east, inner sep=0pt]
    at ($(current page.south east)+(-0.60cm,0.32cm)$)
    {\color{ubHead}\fontsize{7.4}{9}\selectfont #1 / 6
     \ \ \textperiodcentered\ \ \insertframenumber};
\end{tikzpicture}
\end{frame}%
}

% --------------------------------------------------------- peta pertemuan
\newcommand{\blokpeta}[6]{%
  \ifnum#2=#3
    \node[rounded corners=3pt, fill=ubNavy, draw=ubNavy, line width=0.5pt,
          text width=3.02cm, minimum height=1.62cm, align=left, inner sep=5pt]
      at (#1,0) {%
        {\color{ubAccent}\bfseries\fontsize{6.8}{8}\selectfont #4}\\[1pt]
        {\color{white}\bfseries\fontsize{7.6}{9}\selectfont #5}\\[2pt]
        {\color{ubHead}\fontsize{6.4}{7.6}\selectfont #6}};
  \else
    \node[rounded corners=3pt, fill=ubSoft, draw=ubHead, line width=0.5pt,
          text width=3.02cm, minimum height=1.62cm, align=left, inner sep=5pt]
      at (#1,0) {%
        {\color{ubAccent}\bfseries\fontsize{6.8}{8}\selectfont #4}\\[1pt]
        {\color{ubNavy}\bfseries\fontsize{7.6}{9}\selectfont #5}\\[2pt]
        {\color{ubMuted}\fontsize{6.4}{7.6}\selectfont #6}};
  \fi
}

% \petapertemuan{blok aktif}{sebelumnya}{sesudahnya}{keluaran}
% \petapertemuan{blok aktif}{sebelumnya}{pertemuan ini}{sesudahnya}{keluaran}
\newcommand{\petapertemuan}[5]{%
\begin{frame}{Peta Pertemuan}
{\scriptsize\color{ubMuted}Kedudukan materi dalam empat blok perkuliahan}\par
\vspace{0.10cm}
\centering
\begin{tikzpicture}
  \blokpeta{-5.58}{1}{#1}{BLOK I}{Materi 1--5}{Game loop, vektor,\\fisika, collision}
  \blokpeta{-1.86}{2}{#1}{BLOK II}{Materi 6--9}{AI NPC: chase, FSM,\\pathfinding A*}
  \blokpeta{1.86}{3}{#1}{BLOK III}{Materi 10--11}{Unity: scene,\\level, physics}
  \blokpeta{5.58}{4}{#1}{BLOK IV}{Materi 12--14}{Aset visual, audio,\\UI, playtest}
\end{tikzpicture}
\par\vspace{0.30cm}
\begin{columns}[T,onlytextwidth]
  \begin{column}{0.315\textwidth}
    \begin{beamercolorbox}[wd=\linewidth,sep=4pt,rounded=true]{block body}
      {\color{ubMuted}\fontsize{6.6}{8}\selectfont SEBELUMNYA}\\[2pt]
      {\footnotesize #2}
    \end{beamercolorbox}
  \end{column}
  \begin{column}{0.315\textwidth}
    \begin{beamercolorbox}[wd=\linewidth,sep=4pt,rounded=true]{block title}
      {\color{ubHead}\fontsize{6.6}{8}\selectfont PERTEMUAN INI}\\[2pt]
      {\footnotesize #3}
    \end{beamercolorbox}
  \end{column}
  \begin{column}{0.315\textwidth}
    \begin{beamercolorbox}[wd=\linewidth,sep=4pt,rounded=true]{block body}
      {\color{ubMuted}\fontsize{6.6}{8}\selectfont SESUDAHNYA}\\[2pt]
      {\footnotesize #4}
    \end{beamercolorbox}
  \end{column}
\end{columns}
\par\vspace{0.26cm}
{\footnotesize\color{ubText}\textbf{Keluaran pertemuan.} #5\par}
\end{frame}%
}

% --------------------------------------------------------- penutup
\newcommand{\penutup}[2]{%
\begin{frame}[plain]
\begin{tikzpicture}[remember picture, overlay]
  \fill[ubNavy] (current page.south west) rectangle (current page.north east);
  \fill[ubAccent] (current page.north west) rectangle ([xshift=0.10cm]current page.south west);

  \node[anchor=north west, text width=13.6cm, align=left, inner sep=0pt]
    at ($(current page.north west)+(1.00cm,-0.95cm)$) {%
      {\color{ubAccent}\bfseries\fontsize{8.2}{10}\selectfont ASESMEN MINGGU INI}\\[6pt]
      {\color{white}\bfseries\fontsize{24}{28}\selectfont #1}%
    };

  \node[anchor=north west, text width=13.6cm, align=left, inner sep=0pt]
    at ($(current page.north west)+(1.00cm,-4.55cm)$) {%
      {\color{ubNavy!0!ubHead}\rule{13.4cm}{0.4pt}}\\[7pt]
      {\color{ubAccent}\bfseries\fontsize{7.4}{9}\selectfont
       PENGAIT KE MATERI BERIKUTNYA}\\[4pt]
      {\color{ubHead}\fontsize{10.6}{13.2}\selectfont #2}%
    };

  \node[anchor=south east, inner sep=0pt]
    at ($(current page.south east)+(-0.60cm,0.32cm)$)
    {\color{ubHead}\fontsize{7.4}{9}\selectfont \insertframenumber};
\end{tikzpicture}
\end{frame}%
}



\renewcommand{\materino}{07}

\begin{document}

\deckcover{07}{7}
  {Finite State Machine untuk NPC}
  {2.2}
  {Python / Pygame}
  {Bab IX buku ajar}
  {Blok II \textperiodcentered\ Materi 6--9}

% --------------------------------------------------------- capaian
\begin{frame}{Capaian Pertemuan}
\setbeamerfont{itemize/enumerate body}{size=\scriptsize}
\awalan{Sub-CPMK dan indikator dari RPS, serta tujuan Bab IX buku ajar tentang
\emph{Finite State Machine}.}
\begin{columns}[T,onlytextwidth]
  \begin{column}{0.52\textwidth}
    \begin{beamercolorbox}[wd=\linewidth,sep=5pt,rounded=true]{block title}
      {\fontsize{6.8}{8}\selectfont\color{ubHead}SUB-CPMK 2.2}\\[2pt]
      {\footnotesize Mengimplementasikan \emph{Finite State Machine}
       (idle, \emph{patrol}, chase, attack/flee) untuk NPC.}
    \end{beamercolorbox}
    \vspace{0.16cm}

    {\fontsize{7}{8.5}\selectfont\bfseries\color{ubAccent}INDIKATOR}\par
    \vspace{0.06cm}
    \begin{itemize}
      \item Menggambar diagram \emph{state}--transisi NPC.
      \item Mengimplementasikan transisi berbasis jarak dan \emph{cooldown}.
      \item Menguji setiap \emph{state} dan transisinya secara terpisah.
    \end{itemize}
  \end{column}
  \begin{column}{0.45\textwidth}
    {\fontsize{7}{8.5}\selectfont\bfseries\color{ubAccent}TUJUAN BAB IX}\par
    \vspace{0.08cm}
    \begin{itemize}
      \item Menjelaskan mengapa dua keadaan tidak lagi mencukupi.
      \item Menggambar diagram keadaan--transisi dengan syarat yang dapat
            diperiksa.
      \item Memisahkan keputusan dari tindakan di dalam kode.
      \item Menerapkan transisi berdasarkan jarak dan berdasarkan pewaktu.
      \item Mencatat dan memeriksa perpindahan keadaan secara otomatis.
      \item Menjelaskan gejala syarat perpindahan yang kurang lengkap atau
            saling bertentangan.
    \end{itemize}
  \end{column}
\end{columns}
\end{frame}

% --------------------------------------------------------- peta
\petapertemuan{2}
  {Materi 6 \textperiodcentered\ Pengantar AI NPC dan Perilaku Chase}
  {Materi 7 \textperiodcentered\ Finite State Machine untuk NPC}
  {Materi 8 \textperiodcentered\ Pathfinding Grid: dari BFS/Dijkstra ke A*}
  {NPC yang berpatroli, mengejar, dan menyerang menurut diagram keadaan;
   catatan perpindahan dan pemeriksaan keabsahan bersih.}

% ========================================================= bagian 1
\section{Dari Dua Keadaan Menjadi Mesin Keadaan}
\sectionframe{01}{Dari Dua Keadaan Menjadi Mesin Keadaan}
  {Keadaan yang lebih banyak, syarat yang tersendiri, dan penyiapan saat masuk}

\begin{frame}{Dua Keadaan Sudah Merupakan Mesin Keadaan}
\awalan{Variabel \texttt{self.chasing} pada Bab VIII hanya bernilai benar atau
salah. Ia membagi seluruh perilaku NPC menjadi dua golongan: mengejar dan tidak
mengejar.}
\begin{itemize}
  \item Variabel itu sesungguhnya sudah merupakan mesin keadaan dua keadaan.
        Yang belum ada hanyalah nama dan susunannya.
  \item Mesin keadaan yang sesungguhnya menambahkan tiga unsur:
        \begin{itemize}
          \item \textbf{keadaan} (\emph{state}) yang lebih dari dua,
                masing-masing dengan perilaku berbeda;
          \item \textbf{perpindahan} (\emph{transition}) dengan syarat
                tersendiri untuk setiap pasangan keadaan;
          \item \textbf{pekerjaan penyiapan} ketika sebuah keadaan dimasuki.
        \end{itemize}
  \item Setiap perpindahan dipicu \emph{event} yang dapat diperiksa program:
        perbandingan jarak, keadaan tombol, atau pewaktu.
\end{itemize}
\vspace{0.12cm}
\begin{kotak}
NPC yang berpatroli berperilaku berbeda dari NPC yang berdiam, walaupun keduanya
bukan sedang mengejar. Perbedaan itulah yang menuntut keadaan ketiga.
\end{kotak}
\end{frame}

\begin{frame}{Empat Keadaan pada Bab Ini}
\awalan{NPC pada bab ini memiliki empat keadaan, sesuai urutan yang disebut pada
RPS.}
\begin{itemize}
  \item \textbf{berdiam} (\emph{idle}) --- berdiri tidak bergerak sambil
        menghitung lama ia berdiam. Setelah 1,5 detik, ia mulai berpatroli.
  \item \textbf{berpatroli} (\emph{patrol}) --- berpindah dari satu titik
        patroli ke titik berikutnya. Apabila pemain memasuki jari-jari persepsi,
        ia beralih mengejar.
  \item \textbf{mengejar} (\emph{chase}) --- bergerak menuju pemain dengan
        kecepatan tetap. Apabila pemain sudah cukup dekat, ia menyerang.
  \item \textbf{menyerang} (\emph{attack}) --- berhenti dan menyerang secara
        berkala. Apabila pemain menjauh melampaui jangkauan serang, ia kembali
        mengejar.
\end{itemize}
\vspace{0.14cm}
\begin{kotak}
Keempat keadaan dihubungkan oleh tujuh perpindahan yang sah. Tidak ada
perpindahan lain yang diperbolehkan.
\end{kotak}
\end{frame}

\begin{frame}{Mengapa Tidak Cukup Rangkaian Pemeriksaan}
\awalan{Pertanyaan yang wajar: mengapa tidak menuliskan seluruh perilaku sebagai
rangkaian pemeriksaan bersarang, tanpa mesin keadaan sama sekali?}
\begin{itemize}
  \item Cara bersarang berfungsi untuk empat keadaan, tetapi
        \textbf{pekerjaan penyiapan sulit ditempatkan}: menentukan titik
        patroli berikutnya atau mengembalikan pewaktu ke nol harus ditulis di
        dalam cabang, lalu disalin ke cabang lain.
  \item \textbf{Keputusan dan tindakan tercampur}: syarat perpindahan dan
        perilaku keadaan berada di satu tempat, sehingga syarat tidak dapat
        diperiksa tanpa menjalankan perilakunya.
  \item Akibatnya pengujian menjadi sulit, dan cabang-cabang memanjang ketika
        jumlah keadaan bertambah.
  \item Mesin keadaan memisahkan keduanya: syarat perpindahan dikumpulkan di
        satu tempat, perilaku keadaan di tempat lain.
\end{itemize}
\vspace{0.12cm}
\begin{kotak}
Pemisahan itulah yang menjadi pokok dua bagian berikutnya, dan yang membuat
pengujian pada Bagian 4 menjadi mungkin.
\end{kotak}
\end{frame}

% ========================================================= bagian 2
\section{Menggambar Diagram Keadaan--Transisi}
\sectionframe{02}{Menggambar Diagram Keadaan--Transisi}
  {Tiga lambang, tujuh perpindahan, dan tiga jenis pewaktu}

\begin{frame}{Tiga Lambang Diagram Keadaan--Transisi}
\awalan{Diagram digambar sebelum mesin keadaan ditulis sebagai kode. Gambar itu
adalah rancangan, bukan hiasan: kesalahan pada gambar akan menjadi kesalahan
pada program.}
\begin{center}
\adjustbox{max width=0.98\linewidth,max height=0.44\textheight}{%
  % chapters/09/figures/diagram-keadaan.tikz
% Diagram keadaan-transisi NPC: empat keadaan dan tujuh perpindahan yang sah.
% Keterangan pada anak panah dibuat singkat; syarat lengkapnya ada pada tabel
% di dalam teks. Keterangan disusun bertingkat agar tidak saling bertumpuk.
\begin{tikzpicture}[
  keadaan/.style={
    rectangle, rounded corners=4pt, draw=ubPrimary, thick, fill=ubAlt,
    align=center, text width=24mm, inner sep=4pt, minimum height=12mm,
    font=\footnotesize\bfseries
  },
  serang/.style={keadaan, draw=ubRed, fill=ubRed!12},
  panah/.style={-{Stealth[length=2.2mm]}, thick, draw=ubAccent},
  panahbalik/.style={-{Stealth[length=2.2mm]}, thick, draw=ubWarn},
  label/.style={font=\scriptsize, align=center},
  labelbalik/.style={font=\scriptsize, align=center, text=ubWarn}
]

\node[keadaan] (diam) at (0,0) {berdiam};
\node[keadaan] (patrol) at (3.6,0) {berpatroli};
\node[keadaan] (kejar) at (7.2,0) {mengejar};
\node[serang] (serang) at (10.8,0) {menyerang};

% Perpindahan maju, keterangannya disusun bertingkat.
\draw[panah] (diam) -- (patrol);
\node[label, anchor=south] at (1.8,0.78) {1,5 detik};

\draw[panah] (patrol) -- (kejar);
\node[label, anchor=south] at (5.4,1.75) {pemain terlihat};

\draw[panah] (kejar) -- (serang);
\node[label, anchor=south, text width=34mm] at (9.0,0.78)
  {dalam jangkauan serang};

% Perpindahan kembali.
\draw[panahbalik] (serang.south) .. controls (10.8,-1.0) and (7.2,-1.0) ..
  (kejar.south);
\node[labelbalik, anchor=north, text width=34mm] at (9.0,-1.15)
  {keluar dari jangkauan serang};

\draw[panahbalik] (kejar.south) .. controls (7.2,-2.4) and (0,-2.4) ..
  (diam.south);
\node[labelbalik, anchor=north] at (3.6,-2.55) {hilang 1 detik};

\draw[panahbalik] (patrol.north) .. controls (3.6,2.4) and (0,2.4) ..
  (diam.north);
\node[labelbalik, anchor=south] at (1.8,2.55) {putaran selesai};
\end{tikzpicture}
}
\end{center}
\vspace{-0.04cm}
\begin{itemize}
  \item \textbf{Kotak bersudut tumpul} menyatakan sebuah keadaan; isinya nama
        keadaan itu.
  \item \textbf{Anak panah} menyatakan perpindahan yang mungkin terjadi, dari
        keadaan asal menuju keadaan tujuan.
  \item \textbf{Keterangan pada anak panah} menyatakan syarat perpindahan,
        yang harus dapat diperiksa program.
\end{itemize}
\end{frame}

\begin{frame}{Tujuh Perpindahan yang Sah}
\awalan{Diagram tersebut memuat tujuh perpindahan. Tabel berikut memuatnya
beserta syarat masing-masing, sehingga dapat diperiksa satu per satu.}
{\scriptsize
\renewcommand{\arraystretch}{1.12}
\begin{tabularx}{\textwidth}{@{}>{\raggedright\arraybackslash}p{2.35cm}
    >{\raggedright\arraybackslash}p{2.35cm}
    >{\raggedright\arraybackslash}p{1.95cm}
    >{\raggedright\arraybackslash}X@{}}
\toprule
\textbf{Dari} & \textbf{Ke} & \textbf{Jenis syarat} & \textbf{Syarat} \\
\midrule
berdiam & berpatroli & pewaktu & Telah berdiam selama 1,5 detik. \\
\addlinespace
berdiam & mengejar & jarak & Pemain berada di dalam jari-jari persepsi. \\
\addlinespace
berpatroli & berdiam & peristiwa & Satu putaran patroli telah selesai. \\
\addlinespace
berpatroli & mengejar & jarak & Pemain berada di dalam jari-jari persepsi. \\
\addlinespace
mengejar & menyerang & jarak & Pemain berada di dalam jangkauan serang. \\
\addlinespace
mengejar & berdiam & pewaktu & Pemain tidak terlihat selama 1 detik. \\
\addlinespace
menyerang & mengejar & jarak & Pemain keluar dari jangkauan serang. \\
\bottomrule
\end{tabularx}}
\vspace{0.10cm}
{\footnotesize\color{ubMuted}Dua keadaan dapat menuju keadaan yang sama melalui
syarat yang sama, misalnya dari berdiam maupun dari berpatroli menuju mengejar.
Hal itu wajar dan tidak menimbulkan persoalan.}
\end{frame}

\begin{frame}{Pewaktu sebagai Syarat Perpindahan}
\awalan{Tiga dari tujuh perpindahan bersyarat pewaktu, bukan jarak. Ketiganya
dijalankan dengan cara yang sama.}
\begin{itemize}
  \item \textbf{Pewaktu berdiam} menghitung lama NPC berdiam, agar keadaan
        berdiam tidak langsung ditinggalkan pada bingkai berikutnya.
  \item \textbf{Pewaktu tenggang} menghitung lama NPC kehilangan pemain, agar
        ia tidak menyerah hanya karena pemain sekejap berada di balik
        rintangan.
  \item \textbf{Pewaktu jeda} (\emph{cooldown}) menghitung jarak waktu antara
        dua serangan, agar serangan tidak terjadi pada setiap bingkai.
  \item Setiap pewaktu ditambah selang waktu setiap bingkai, lalu dibandingkan
        dengan nilai batasnya.
\end{itemize}
\vspace{0.12cm}
\begin{peringatan}
Tiga kesalahan diagram yang sering terjadi: perpindahan tanpa syarat yang dapat
diperiksa, keadaan tanpa jalan keluar, dan dua syarat yang saling bertentangan.
Gejala terakhir khas: NPC sesekali berperilaku tidak seperti dugaan, dan hanya
pada keadaan tertentu.
\end{peringatan}
\end{frame}

% ========================================================= bagian 3
\section{Menerapkan Mesin Keadaan di Dalam Kode}
\sectionframe{03}{Menerapkan Mesin Keadaan di Dalam Kode}
  {Keputusan, penyiapan, dan tindakan dipisahkan}

\begin{frame}{Tiga Pekerjaan yang Dipisahkan}
\awalan{Susunan program memisahkan tiga pekerjaan yang berbeda. Pemisahan
itulah yang menjadi inti bagian ini.}
\begin{itemize}
  \item \texttt{decide} membaca keadaan sekarang, jarak terhadap pemain, dan
        ketiga pewaktu, lalu mengembalikan nama keadaan berikutnya. Fungsi ini
        tidak mengubah apa pun.
  \item \texttt{enter\_state} menetapkan keadaan baru, mengembalikan penghitung
        keadaan itu ke nol, dan mencatat perpindahannya.
  \item \texttt{act} menjalankan perilaku keadaan yang sedang aktif, yaitu
        mengubah kecepatan NPC.
  \item \texttt{update} menyatukan ketiganya: menghitung jarak, meminta
        keputusan, menyiapkan keadaan baru apabila berpindah, lalu bertindak.
\end{itemize}
\vspace{0.12cm}
\begin{kotak}
Karena \texttt{decide} tidak mengubah apa pun, keputusan dapat diperiksa dengan
memanggilnya berulang kali tanpa menjalankan permainan.
\end{kotak}
\end{frame}

\begin{frame}[fragile]{Syarat Perpindahan Ditulis pada Satu Tempat}
\awalan{Seluruh syarat dari tabel dikumpulkan di dalam \texttt{decide}, disusun
menurut keadaan yang sedang aktif.}
\lstinputlisting[language=Python,firstline=160,lastline=170]
  {../buku_ajar/code/python/ch09_fsm.py}
\kodejelas{Pewaktu jeda serangan dijalankan pada awal fungsi, sebab ia berjalan
tanpa memandang keadaan. Pada keadaan berdiam, syarat ``pemain terlihat''
diperiksa lebih dahulu daripada syarat ``sudah cukup lama berdiam'': NPC yang
melihat pemain harus segera mengejar, walaupun ia baru saja mulai berdiam.}
\end{frame}

\begin{frame}[fragile]{Penghalang pada Batas Jangkauan Serang}
\awalan{Keadaan mengejar memakai jari-jari persepsi dan pewaktu tenggang;
keadaan menyerang memakai penghalang pada batas jangkauan.}
\lstinputlisting[language=Python,firstline=179,lastline=192]
  {../buku_ajar/code/python/ch09_fsm.py}
\kodejelas{Syarat dari menyerang ke mengejar memakai jangkauan serang dikalikan
1,2, bukan jangkauan itu sendiri. Tanpa penghalang, kedua syarat benar bersamaan
pada jarak tepat jangkauan; NPC berganti keadaan setiap bingkai dan tampak
berkedip serta tersendat di tepi jangkauan.}
\end{frame}

% ========================================================= bagian 4
\section{Menguji Setiap Keadaan dan Perpindahannya}
\sectionframe{04}{Menguji Setiap Keadaan dan Perpindahannya}
  {Catatan, skenario berulang, dan pemeriksaan keabsahan}

\begin{frame}{Mengapa Pengujiannya Sulit}
\awalan{Mesin keadaan sulit diperiksa dengan cara biasa, sebab perilakunya
bergantung pada posisi pemain dan pada waktu.}
\begin{itemize}
  \item Memeriksa keadaan menyerang secara manual menuntut pemain mendekat,
        menunggu NPC menyadarinya, lalu menunggu NPC mendekat. Setiap percobaan
        puluhan detik dan hasilnya bergantung pada keterampilan mengendalikan
        pemain.
  \item Tiga cara dipakai bersama: mencatat perpindahan beserta nomor
        bingkainya, menjalankan skenario yang dapat diulang, dan memeriksa
        keabsahan setiap perpindahan.
  \item Catatan perpindahan lebih berguna daripada tangkapan layar: ia
        memperlihatkan urutan dan lama setiap keadaan, bukan hanya keadaan yang
        sedang aktif.
\end{itemize}
\vspace{0.12cm}
\begin{kotak}
Skenario menuntut pemain berdiri cukup jauh ketika menghilang, kira-kira 410
piksel pada nilai bab ini. Bila kurang jauh, NPC menemukan pemain kembali
sebelum pewaktu tenggangnya habis.
\end{kotak}
\end{frame}

\begin{frame}[fragile]{Memeriksa Keabsahan Perpindahan}
\awalan{Tabel perpindahan yang sah adalah diagram keadaan yang ditulis kembali
dalam bentuk yang dapat diperiksa program.}
\lstinputlisting[language=Python,firstline=58,lastline=63]
  {../buku_ajar/code/python/ch09_fsm.py}
\kodejelas{Setiap pasangan keadaan yang berurutan pada catatan perpindahan
diperiksa terhadap tabel ini. Pemeriksaan sifat yang harus selalu berlaku tidak
menuntut urutan tertentu, sehingga tahan terhadap perbedaan kecepatan komputer,
tetapi tetap menangkap perpindahan yang keliru.}
\end{frame}

\begin{frame}{Memandang Keadaan sebagai Garis Waktu}
\awalan{Catatan perpindahan dapat digambarkan sebagai garis waktu, sehingga lama
setiap keadaan terlihat sekaligus.}
\begin{center}
\adjustbox{max width=0.99\linewidth,max height=0.30\textheight}{%
  % chapters/09/figures/garis-waktu-keadaan.tikz
% Garis waktu keadaan NPC pada skenario pengujian, memakai catatan perpindahan
% yang sesungguhnya: setiap pita menyatakan lama sebuah keadaan berlangsung.
\begin{tikzpicture}[
  pita/.style={rectangle, draw=ubLine, thick, minimum height=9mm,
               align=center, font=\scriptsize},
  label/.style={font=\scriptsize, align=center},
  penanda/.style={draw=ubLine, thin}
]

% Pita keadaan. Lebar sebanding dengan lama keadaan berlangsung.
\node[pita, fill=ubGray, minimum width=1.42cm] (a) at (0.71,0) {berdiam};
\node[pita, fill=ubGreenBg, minimum width=0.93cm, right=0cm of a] (b) {patroli};
\node[pita, fill=ubRedBg, minimum width=1.03cm, right=0cm of b] (c) {kejar};
\node[pita, fill=ubWarnBg, minimum width=3.17cm, right=0cm of c] (d) {menyerang};
\node[pita, fill=ubRedBg, minimum width=0.95cm, right=0cm of d] (e) {kejar};
\node[pita, fill=ubGray, minimum width=1.42cm, right=0cm of e] (f) {berdiam};
\node[pita, fill=ubGreenBg, minimum width=5.09cm, right=0cm of f] (g) {berpatroli};

% Nomor bingkai pada setiap batas keadaan.
\foreach \titik/\bingkai in {a.west/0, a.east/91, b.east/151, c.east/217,
                             d.east/421, e.east/482, f.east/573, g.east/900} {
  \draw[penanda] (\titik) ++(0,0.75) -- ++(0,0.25);
  \node[label, anchor=south] at ($(\titik) + (0,1.0)$) {\bingkai};
}

% Sumbu waktu.
\draw[-{Stealth[length=2mm]}, thick, draw=ubLine] (-0.3,-0.9) -- (14.6,-0.9);
\node[label, anchor=north west] at (0,-1.1)
  {bingkai (900 bingkai setara 15 detik pada 60 bingkai per detik)};

\node[label, anchor=north west, text width=140mm] at (-0.3,-1.9)
  {Pita berwarna menyatakan keadaan yang sedang aktif, dan lebarnya sebanding
   dengan lama keadaan itu berlangsung. Perhatikan bahwa keadaan menyerang
   berlangsung paling lama, yaitu 204 bingkai, sebab pemain berdiri di dalam
   jangkauan serang selama 3,4 detik.};
\end{tikzpicture}
}
\end{center}
\vspace{-0.06cm}
\begin{itemize}
  \item Skenario mencatat urutan: 0 berdiam, 91 berpatroli, 151 mengejar,
        217 menyerang, 421 mengejar, 482 berdiam, 571 berpatroli.
  \item Keadaan menyerang berlangsung paling lama, 204 bingkai atau 3,4 detik.
        Dengan pewaktu jeda 0,8 detik, jumlah serangan seharusnya
        $3{,}4 \div 0{,}8 + 1 \approx 5$, dan program mencatat lima serangan.
  \item Apabila hasil tidak sesuai dugaan, periksa lebih dahulu apakah dugaan
        atau rancangan skenarionya yang keliru, bukan programnya.
\end{itemize}
\end{frame}

% ========================================================= bagian 5
\section{Program Mesin Keadaan Lengkap}
\sectionframe{05}{Program Mesin Keadaan Lengkap}
  {Susunan kelas, keluaran pengujian, dan percobaan yang disarankan}

\begin{frame}{Susunan Kelas \emph{Npc}}
\awalan{Seluruh perilaku NPC terletak pada kelas \texttt{Npc}. Fungsi
\texttt{update} hanya berisi empat langkah dan tidak memuat satu pun syarat
perpindahan.}
\begin{itemize}
  \item \texttt{decide} menentukan keadaan berikutnya; ia hanya membaca dan
        mengembalikan jawaban.
  \item \texttt{enter\_state} menyiapkan keadaan baru dan mencatat
        perpindahannya.
  \item \texttt{act} menjalankan perilaku keadaan yang sedang aktif.
  \item \texttt{update} menghitung jarak, meminta keputusan, menyiapkan keadaan
        baru apabila berpindah, lalu bertindak.
\end{itemize}
\vspace{0.14cm}
\begin{kotak}
Seluruh syarat perpindahan berada di dalam \texttt{decide}. Susunan itulah yang
membuat keputusan dapat diperiksa terpisah dari perilakunya.
\end{kotak}
\end{frame}

\begin{frame}[fragile]{Keluaran Pengujian}
\awalan{Dijalankan dengan skenario, program menampilkan catatan perpindahan
beserta hasil pemeriksaannya.}
\begin{lstlisting}[language={},numbers=none]
Bingkai dijalankan : 900
Serangan dilakukan : 5
Catatan perpindahan keadaan (bingkai, keadaan):
        0  berdiam
       91  berpatroli
      151  mengejar
      217  menyerang
      421  mengejar
      482  berdiam
      571  berpatroli
Hasil pemeriksaan  : seluruh perpindahan sah
Hasil skenario     : SESUAI dengan urutan yang diharapkan
\end{lstlisting}
\kodejelas{Ketujuh perpindahan terjadi seluruhnya, dalam urutan yang dapat
diperkirakan, dan seluruhnya sah menurut tabel perpindahan. Empat keadaan pada
diagram telah diuji satu per satu.}
\end{frame}

\begin{frame}{Percobaan yang Disarankan}
\awalan{Mengubah satu nilai pada satu waktu memperlihatkan tugas masing-masing
unsur mesin keadaan.}
{\scriptsize
\renewcommand{\arraystretch}{1.12}
\begin{tabularx}{\textwidth}{@{}>{\raggedright\arraybackslash}p{6.0cm}
    >{\raggedright\arraybackslash}X@{}}
\toprule
\textbf{Yang diubah dan caranya} & \textbf{Yang sebaiknya diamati} \\
\midrule
\texttt{ATTACK\_COOLDOWN} menjadi 0,2 detik &
Jumlah serangan bertambah, sedangkan urutan keadaannya tidak berubah. \\
\addlinespace
\texttt{GIVE\_UP\_SECONDS} menjadi 3 detik &
NPC mengejar lebih lama; pemain harus berdiri lebih jauh agar benar-benar
kehilangan jejak. \\
\addlinespace
\texttt{IDLE\_SECONDS} menjadi nol &
Keadaan berdiam tidak lagi berlangsung lama; NPC langsung berpatroli sesudah
menyerah. \\
\addlinespace
Hapus perkalian 1,2 pada jangkauan serang &
NPC berkedip dan tersendat pada tepi jangkauan serang, sebab kedua syarat benar
bersamaan. \\
\addlinespace
Tambahkan perpindahan yang tidak ada pada diagram &
Pemeriksaan keabsahan melaporkan perpindahan yang tidak sah. \\
\bottomrule
\end{tabularx}}
\vspace{0.10cm}
\begin{kotak}
Percobaan terakhir paling penting. Menambahkan satu perpindahan yang tidak ada
pada diagram akan langsung ditangkap pemeriksaan keabsahan; diagram menjadi
ukuran kebenaran program, bukan sekadar catatan.
\end{kotak}
\end{frame}

% ========================================================= bagian 6
\section{Yang Belum Ada pada Mesin Keadaan Ini}
\sectionframe{06}{Yang Belum Ada pada Mesin Keadaan Ini}
  {Niat tanpa rute, dan bentuk mesin keadaan yang lebih besar}

\begin{frame}{Keputusan tanpa Pengetahuan tentang Jalan}
\awalan{Mesin keadaan menjawab pertanyaan ``apa yang sebaiknya saya lakukan?'',
bukan ``bagaimana cara menuju ke sana?''.}
\begin{itemize}
  \item NPC memutuskan untuk mengejar, tetapi pelaksanaannya hanya satu:
        bergerak lurus ke arah pemain.
  \item Empat keadaan tidak menolong NPC yang terhalang dinding; keadaan
        mengejar tetap menekan dinding selama pemain berada di baliknya.
  \item Yang bertambah hanyalah keragaman perilaku, bukan keterangan tentang
        jalan.
  \item Bab X menambahkan pencarian jalur: keadaan mengejar tidak lagi
        menghasilkan gerak lurus, melainkan memanggil A* untuk menentukan rute.
\end{itemize}
\vspace{0.12cm}
\begin{kotak}
Mesin keadaan menentukan \textbf{niat}, sedangkan pencarian jalur menentukan
\textbf{rute}. Pemisahan tanggung jawab itu mulai terbentuk di sini.
\end{kotak}
\end{frame}

\begin{frame}{Bentuk Mesin Keadaan yang Lebih Besar}
\awalan{Dua bentuk pengembangan lazim pada permainan yang lebih besar; keduanya
bertolak dari gagasan yang sama dengan mesin keadaan.}
\begin{itemize}
  \item \textbf{Mesin keadaan bertingkat} (\emph{hierarchical state machine}):
        sebuah keadaan dapat memuat mesin keadaan yang lebih kecil. Misalnya,
        keadaan mengejar memuat keadaan menghindari rintangan dan mendekati
        pemain.
  \item \textbf{Pohon perilaku} (\emph{behaviour tree}): perilaku disusun
        sebagai pohon keputusan yang dibaca dari atas ke bawah; setiap cabang
        menyatakan satu pertimbangan.
  \item Pohon perilaku lebih mudah dibaca ketika jumlah perilakunya banyak, dan
        banyak dipakai pada permainan yang sesungguhnya.
  \item Keduanya di luar cakupan mata kuliah, tetapi sama-sama memisahkan
        keadaan yang mungkin terjadi dari syarat perpindahannya.
\end{itemize}
\vspace{0.12cm}
\begin{kotak}
Cara itu mencegah jumlah keadaan bertambah secara berlipat ketika perilaku NPC
makin beragam.
\end{kotak}
\end{frame}

\begin{frame}{Ringkasan Perkembangan}
\awalan{Bab ini melengkapi bagian pertama lapisan keputusan; yang masih kurang
adalah kemampuan NPC mengetahui jalannya sendiri.}
{\scriptsize
\renewcommand{\arraystretch}{1.12}
\begin{tabularx}{\textwidth}{@{}>{\raggedright\arraybackslash}p{1.7cm}
    >{\raggedright\arraybackslash}p{4.9cm}
    >{\raggedright\arraybackslash}X@{}}
\toprule
\textbf{Bab} & \textbf{Yang ditambahkan} & \textbf{Akibatnya} \\
\midrule
Bab X & Pencarian jalur A* pada peta berhalangan &
Keadaan mengejar tidak lagi menghasilkan gerak lurus, melainkan rute yang
menghindari rintangan. \\
\addlinespace
Bab XI & Penggabungan mesin keadaan dan pencarian jalur &
NPC mengejar dengan rute memutar, dan biaya perhitungannya diatur agar tetap
ringan. \\
\addlinespace
Bab XIII & Fisika mesin pada Unity &
Sebagian perilaku dapat disusun memakai perkakas mesin, sedangkan logika
keputusannya tetap ditulis sendiri. \\
\bottomrule
\end{tabularx}}
\vspace{0.10cm}
\begin{kotak}
Setelah bab ini, seluruh perkakas untuk menyusun perilaku NPC yang beragam sudah
tersedia. Yang masih kurang hanyalah kemampuan NPC untuk mengetahui jalannya
sendiri.
\end{kotak}
\end{frame}

% --------------------------------------------------------- aktivitas
\begin{frame}{Aktivitas di Kelas \textperiodcentered\ Praktikum 07}
\awalan{Pengalaman belajar minggu 7 pada RPS, dikerjakan berurutan di kelas lalu
diselesaikan di rumah.}
\begin{enumerate}
  \item Ketik ulang \texttt{ch09\_fsm.py} ke folder \texttt{code}, jalankan,
        lalu amati warna NPC yang berubah menurut keadaannya dan jalur patroli
        yang dilaluinya.
  \item Amati setiap keadaan secara terpisah: berdiamlah di jauh sampai satu
        putaran patroli selesai, dekati sampai NPC mengejar, lalu berdirilah di
        dekatnya sampai ia menyerang. Catat keadaan yang paling lama
        berlangsung.
  \item Jalankan skenario pengujian
        \texttt{python ch09\_fsm.py --frames 900 --scenario}, lalu catat
        seluruh perpindahan beserta nomor bingkainya.
  \item Gambarlah diagram keadaan Anda sendiri dengan empat keadaan dan syarat
        pada setiap perpindahan; bandingkan dengan diagram bab, lalu periksa
        apakah setiap syarat benar-benar ada di dalam kode.
  \item Tambahkan keadaan menghindar ketika nyawa tinggal sedikit, tentukan
        syarat masuk dan keluarnya, lalu uji satu perpindahan yang tidak ada
        pada diagram dan amati laporan pemeriksaan keabsahan.
  \item Tulis catatan praktikum yang memuat diagram, catatan perpindahan, dan
        hasil pemeriksaan keabsahan; simpan pekerjaan ke repositori Git.
\end{enumerate}
\vspace{0.10cm}
{\footnotesize\color{ubAccent}\itshape Rincian lengkap: Praktikum 07 pada Bab IX
buku ajar. Kumpulkan melalui repositori Git pribadi.}
\end{frame}

% --------------------------------------------------------- rangkuman
\begin{frame}{Rangkuman}
\begin{itemize}
  \item Variabel bernilai benar-salah sudah merupakan mesin keadaan dua
        keadaan; mesin yang sesungguhnya menambahkan keadaan lebih dari dua,
        syarat perpindahan tersendiri, dan penyiapan saat masuk keadaan.
  \item NPC memiliki empat keadaan: berdiam, berpatroli, mengejar, dan
        menyerang, dengan tujuh perpindahan yang sah.
  \item Diagram memakai tiga lambang; setiap perpindahan harus memiliki syarat
        yang dapat diperiksa, setiap keadaan harus memiliki jalan keluar, dan
        tidak boleh ada dua syarat yang saling bertentangan.
  \item Tiga pewaktu dipakai: pewaktu berdiam, pewaktu tenggang, dan pewaktu
        jeda serangan (\emph{cooldown}); ketiganya ditambah selang waktu lalu
        dibandingkan dengan nilai batasnya.
  \item Kode memisahkan \texttt{decide}, \texttt{enter\_state}, dan
        \texttt{act}; \texttt{decide} tidak mengubah apa pun sehingga keputusan
        dapat diuji terpisah dari tindakan.
  \item Pengujian memakai catatan perpindahan beserta bingkainya, skenario yang
        dapat diulang, dan pemeriksaan keabsahan; mesin keadaan menentukan niat,
        bukan rute.
\end{itemize}
\end{frame}

% --------------------------------------------------------- cek
\begin{frame}{Cek Pemahaman}
\begin{enumerate}
  \item Sebutkan tiga unsur yang ditambahkan mesin keadaan dibandingkan variabel
        bernilai benar-salah, dan jelaskan mengapa masing-masing diperlukan.
  \item Mengapa fungsi \texttt{decide} sengaja tidak mengubah keadaan apa pun,
        dan apa manfaatnya bagi pengujian?
  \item Mengapa pemeriksaan keabsahan perpindahan lebih berguna daripada
        pemeriksaan urutan keadaan?
\end{enumerate}
\vspace{0.35cm}
{\footnotesize\color{ubMuted}\itshape Jawab tanpa membuka catatan, lalu
bandingkan dengan pasangan di sebelah Anda.}
\end{frame}

% --------------------------------------------------------- penutup
\penutup{Tugas Praktikum 07}
  {Materi 8 membuka pathfinding grid: keadaan mengejar tidak lagi bergerak
   lurus, melainkan memanggil pencarian jalur untuk menentukan rute. Anda akan
   membangun grid berbiaya, memakai \emph{heapq}, dan merekonstruksi jalur.}

\end{document}
